\chapter{The Officer, Again}
\label{sec:35}
We started with an officer who had twenty seconds per name on machine-generated lists of people to potentially kill by airstrike. He was a component in a sociotechnical pipeline, and he was also operating that pipeline. He had signed up for this role in direct response to the October 7, 2023 attacks, when Hamas killed 1,200 Israelis — most of them noncombatants. He and his colleagues were pressured from all sides, including inside, to approve the names on the lists. They were reeling with grief and vengefulness.

I have found it challenging to convey to friends and relatives the depth of my disappointment in Israeli society, my concern for the threats to personal safety that Israel has been creating for diasporic Jews, and my belief that Israel, insofar as it is an administrative state, has forfeited its right to exist.

Since October 7, the Israel Defense Forces have killed 100,000 Palestinians \autocite{gomezugarte2025uncertainty,acosta2026gaza,economist2025gaza}. While writing, three years later, I asked ChatGPT: what could Israel have done differently during the first two months of the conflict (October–November 2023) that may have reduced noncombatant fatalities? The reply began: “Israel could likely have reduced noncombatant deaths substantially by limiting bombardment in populated areas, maintaining essential supplies and medical care, making evacuations safer, and pursuing an earlier or longer negotiated halt to fighting.” Gaza had one of the highest population densities in the world. It does not take an advanced supercomputer to foresee the outcomes of using heavy explosive weapons in these neighborhoods. There may have been a justifiable military counterattack to October 7 (personally I am not sure), but what I am certain of is that what Israel has \emph{actually} done since that day is indefensible.

(Hamas has also forfeited its right to exist. As has the Roberts Court. In none of these three cases is heavy aerial bombardment an appropriate remedy.)

What was it that Rabbi Hillel said when asked to teach the whole Torah while standing on one foot?

Richard Wright admired H. L. Mencken for how he used words as weapons. I hope that going forward, my fellow Jews (who sometimes refer to themselves, in true model-minority fashion, as “people of the book”), particularly the Zionist ones, will refamiliarize themselves with the cost-effectiveness and safety of using words rather than genocide, and particularly in the reading of words rather than in the publishing of them.

Was it about Hashem? Did Hillel say something important about Hashem?

Zionism is destroying Judaism. By entangling itself with Judaism, it has turned Judaism itself into an expression of settler-colonialism. This is not something I will forgive within the foreseeable future. People can be really damn disappointing.

If I have learned anything from watching \emph{Neon Genesis Evangelion}, it is that the inevitability of people disappointing you is no reason to give up on people generally. Similarly, just because the Israeli state and political order has failed does not mean people should give up on politics generally. (Not that they could even if they wanted to. But what I mean is that people should not give up on assassination-free politics, abduction-free politics, warcrime-free politics.)

Writing in this way is a combination of catharsis, bitterness, triumphalism, vengefulness, and certitude. But if I do not hold Israel accountable, then who will?

Five years ago I would have said that Zionism is not inherently racist, and that in spite of political setbacks, the concept of Zionism was spacious enough for a binational one-state solution: two peoples on one strip of land, each with the same standing, neither one the other's object.

Whatever the word Zionism might have meant five years ago, today's Zionism is smaller. One people holds a state, an army, an air force and the crossings, and the other people lives under them. A commission of inquiry has found that Israel is committing genocide in Gaza \autocite{coiopt2025genocide}. The International Court of Justice has found the occupation unlawful and said it must end \autocite{icj2024palestine}.

When one side holds every lever, equal standing doesn't come from being evenhanded in your slogans. It starts with the side that doesn't have it. That's what Free Palestine means: Palestinians free of the siege, the occupation and the bombs, with the standing every Israeli already has. It isn't about driving anyone out of anywhere. And also, Black lives matter. That doesn't take away from White lives, Brown lives, any other kind of lives mattering. It is simply an affirmation that Black lives matter. There is not a finite amount of mattering. The supply of mattering expands endogenously with the lives. There is enough for everyone.

Somewhere between Atlas shrugging and eugenics, there ought to be room for fallible people helping one another develop healthier relationships to their feelings, and taking responsibility for what they do with them. There is no code for the human spirit. And people are not products. But they do reproduce. As do their ideas.

Was it about Amalek? Was it about Jericho? Is that how Hillel answered that wicked child, forcing a man of his age to do gymnastics?

The work I have in mind involves learning how guilt can prompt repair, compassion can widen attention, and grief can be endured without becoming a warrant for someone else's suffering. The work can give us practice in noticing harm, revising our judgments, and refusing instructions.

Oh, now I remember what Hillel said. It was “Free Palestine. The rest is just commentary.”
